% Generated draft: requires mathematical and source review.
\documentclass[11pt]{article}
\usepackage[margin=25mm]{geometry}
\usepackage{fontspec}
\setmainfont{Latin Modern Roman}
\usepackage{amsmath,amssymb,mathtools,unicode-math}
\setmathfont{Latin Modern Math}
\usepackage{enumitem,needspace,xurl,hyperref}
\setlength{\parindent}{0pt}
\setlength{\parskip}{5pt}
\setlength{\emergencystretch}{3em}
\newcommand{\N}{\mathbb{N}}
\newcommand{\Z}{\mathbb{Z}}
\newcommand{\Q}{\mathbb{Q}}
\newcommand{\R}{\mathbb{R}}
\newcommand{\C}{\mathbb{C}}
\newcommand{\F}{\mathbb{F}}
\newcommand{\A}{\mathbb{A}}
\newcommand{\PP}{\mathbb{P}}
\newcommand{\E}{\mathbb{E}}
\newcommand{\eps}{\varepsilon}
\newcommand{\dd}{\,\mathrm{d}}
\begin{document}
\section*{UnsolvedMath ID: 700473}
\textbf{Source problem code:} catalogue entry
\subsection{Definitions and mathematical statement}
A variety of groups is a class of groups defined by a collection of group identities. A group \(G\) is relatively free if, for some variety \(\mathcal V\), it belongs to \(\mathcal V\) and has a generating set \(X\) such that every map from \(X\) to any group \(H\in\mathcal V\) extends uniquely to a homomorphism \(G\to H\). No restriction on the cardinality of \(X\) is imposed in the supplied statement.

Define the commutator subgroup by
\[
G'=\langle [a,b]:a,b\in G\rangle,
\qquad [a,b]=a^{-1}b^{-1}ab.
\]
A group is periodic if each of its elements has finite order. It has finite exponent if there is a positive integer that annihilates every element by taking that power.

The question is whether every relatively free group \(G\) satisfies
\[
\left(\forall h\in G'\;\exists m\in\Z_{>0}:h^m=1\right)
\quad\Longrightarrow\quad
\left(\exists n\in\Z_{>0}\;\forall h\in G':h^n=1\right).
\]
Here \(n\) may depend on \(G\); the conclusion requires a single bound for all elements of its commutator subgroup.
\subsection{Short English statement}
If every element of the commutator subgroup of a relatively free group has finite order, must there be one positive integer \(n\) such that every element of that subgroup has order dividing \(n\)?
\subsection{Sources}
\begin{itemize}
\item Local UnsolvedMath registry, supplied complete source record: permanent ID \(700473\), exact source code \texttt{catalogue entry}, title ``Kourovka Notebook Problem 8.21.''
\item \emph{The Kourovka Notebook (Unsolved Problems in Group Theory)}, original problem 8.21, 8th issue (1982); proposed by L. G. Kov\'acs. The supplied record cites the 21st issue edition (2026), \texttt{arXiv:1401.0300v48}, and identifies E. I. Khukhro and V. D. Mazurov as editors of the Issues 1--21 collection.
\item The registry labels this problem ``open.'' This is a supplied source classification, not an independently verified assessment of its current mathematical status. The cited source has not been independently consulted for this draft.
\end{itemize}
\end{document}
